
\subsubsection{Initialization Symmetry (H-E1)}

\begin{table}[htbp]
\centering
\resizebox{\columnwidth}{!}{%
\begin{tabular}{ll}
\toprule
Seed & $SR_0$ \\
\midrule
0 & 1.0028 \\
1 & 0.9948 \\
2 & 0.9974 \\
3 & 1.0037 \\
4 & 1.0008 \\
\bottomrule
\end{tabular}
}
\end{table}

Mean $SR_0$ = 0.9999, 95\% CI = [0.9953, 1.0046], SEM = 0.0017.

Both success criteria were satisfied: (1) mean $SR_0 \in$ [0.9, 1.1], (2) 95\% CI includes 1.0.

\subsubsection{Temporal Precedence (H-M1)}

\begin{table}[htbp]
\centering
\resizebox{\columnwidth}{!}{%
\begin{tabular}{ll}
\toprule
Seed & $\tau _{r}\rightarrow SR$ \\
\midrule
0 & 4 \\
1 & 3 \\
2 & 4 \\
\bottomrule
\end{tabular}
}
\end{table}

Mean $\tau$ = 3.67 epochs, 95\% CI = [3.0, 4.0]. The confidence interval excludes zero.

Observed trajectories: gradient ratio r\_t decayed from approximately 1.02 to approximately 0.65 over 10 epochs; sharpness ratio SR\_t increased from approximately 1.01 to approximately 1.72 over 10 epochs.

\subsubsection{Update-Norm Parity Intervention (H-M2)}

Implementation status: Complete (10 code files, approximately 600 lines).
Code validation: All modules pass import tests, per-group gradient norm computation verified, parity scaling correctly interpolates toward mean norm.
Experiment execution: Not completed due to CPU-only hardware constraints.

Unit test observations: group gradient norms varied from 5.27 to 9.63; parity scaling reduced this variance (target norm = 7.35).

\subsubsection{Summary}

\begin{table}[htbp]
\centering
\resizebox{\columnwidth}{!}{%
\begin{tabular}{lll}
\toprule
Prediction & Result & Status \\
\midrule
P1: $\tau _{r}\rightarrow SR$ > 0 & 3.67 epochs, CI excludes zero & SUPPORTED \\
P3: $SR_0 \approx$ 1.0 & 0.9999, CI includes 1.0 & SUPPORTED \\
P2: Parity attenuates SR & Code validated only & INCONCLUSIVE \\
P4: SR reduction improves WGA & Not tested & NOT TESTED \\
\bottomrule
\end{tabular}
}
\end{table}

